\documentclass[11pt]{article}
\usepackage[a4paper,margin=27mm]{geometry}
\usepackage{amsmath,amssymb,amsthm,mathtools}
\usepackage[colorlinks=true,urlcolor=blue,linkcolor=blue,citecolor=blue]{hyperref}
\usepackage{microtype}
\newtheorem{theorem}{Theorem}
\newtheorem{proposition}[theorem]{Proposition}
\newtheorem{corollary}[theorem]{Corollary}
\newtheorem{remark}[theorem]{Remark}
\DeclareMathOperator{\Ric}{Ric}
\DeclareMathOperator{\Vol}{Vol}
\DeclareMathOperator{\im}{im}
\DeclareMathOperator{\Ext}{Ext}
\newcommand{\Z}{\mathbb Z}
\newcommand{\R}{\mathbb R}
\title{Scalar Curvature, Ambient Fundamental Groups,\ and the Attainment Problem for Curvature Energy}
\author{Research continuation}
\date{12 October 2026}
\begin{document}
\maketitle
\begin{abstract}
We derive a global restriction on compact domains in closed aspherical
four-manifolds from positive scalar curvature and mapping degree. A domain
with finite boundary fundamental group that admits nonnegative scalar
curvature and strictly mean-convex boundary has nullhomotopic inclusion.
Consequently the previous topology and energy restrictions for embedded
ALE cores no longer require a finite fundamental group of the core itself.
For a nontrivial spherical end the resulting trace-free Ricci energy bound
is $4\pi^2+\frac1{12}\int s^2$, with a known scalar-flat quotient realizing
the constant. We also give a complete cusp-doubling calculation explaining
why eliminating Ricci-flat bubbles does not ensure smooth attainment.
These are consequences of established global theorems and explicit
calculations; no priority claim is made. A smooth non-Bach-flat unrestricted
global minimizer has not been constructed or proved to exist here.
\end{abstract}

\section{The variational target}
For a closed smooth four-manifold $M$ write
\[
 P(g)=\int_M|\Ric_g^{\circ}|^2\,dV_g,
 \qquad P_*(M)=\inf_{g\in\operatorname{Met}(M)}P(g).
\]
The target is smooth attainment of this unrestricted infimum by a metric
which is not Bach-flat. With curvature-operator norms,
\[
 \int_M|\operatorname{Rm}_{\mathrm{op}}|^2\,dV_g
 =8\pi^2\chi(M)+P(g).
\]
Thus this is precisely the corresponding full-curvature optimization
question in LeBrun's discussion \cite[Section 11]{LeBrun04}.
The original problem requires both global minimality and smooth attainment.
Neither follows from the critical-point equation alone.

\section{A global scalar-curvature obstruction for embedded domains}
Throughout, aspherical means that the universal cover is contractible.
Mean curvature is computed with the outward unit normal and is positive
for an ordinary Euclidean ball.

\begin{theorem}\label{thm:domain}
Let $M$ be a closed connected oriented smooth aspherical four-manifold,
and let $N_0\subset M$ be a connected compact smooth codimension-zero
submanifold with nonempty connected boundary $Y$.
Suppose that $\pi_1(Y)$ is finite and that $N_0$ admits a smooth metric with
\[
 s\ge0,\qquad H_Y>0.
\]
Then the inclusion $N_0\hookrightarrow M$ is nullhomotopic.
No assumption is made on $\pi_1(N_0)$.
\end{theorem}

\begin{proof}
First the double $DN_0=N_0\cup_Y(-N_0)$ admits a positive-scalar-curvature
metric. Here is the reduction from nonnegative to positive scalar
curvature before applying the classical doubling theorem.
Set $\Delta=\operatorname{div}\nabla$, and solve
\[
 -\Delta w=1\ \text{on }N_0,\qquad w|_Y=0.
\]
For sufficiently small $\varepsilon>0$, $u=1+\varepsilon w>0$ and
\[
 s_{u^2g}=u^{-3}(-6\Delta u+s_gu)
          =u^{-3}(6\varepsilon+s_gu)>0.
\]
Strict mean convexity persists under this small smooth perturbation;
explicitly, when $H$ is the trace of the second fundamental form,
\[
 H_{u^2g}=u^{-1}\bigl(H_g+3u^{-1}\partial_\nu u\bigr).
\]
Gromov--Lawson doubling now gives a smooth positive-scalar-curvature
metric on $DN_0$; the implication needed here is recalled in
\cite[Introduction]{Frenck}. The converse doubling conjecture is not used.

Let $C=\overline{M\setminus N_0}$, which is connected because $Y$ is
connected. Put $G=\pi_1(M)$. This group is torsion-free and one-ended.
Indeed, it is the group of a finite aspherical closed manifold; its
universal cover $U$ is a contractible oriented four-manifold, and
$H_c^1(U;\Z)=H_3(U;\Z)=0$ gives one-endedness. In particular $G$ cannot
be a nontrivial free product.

The homomorphism $\pi_1(Y)\to G$ is trivial. Van Kampen therefore gives
\[
 G\cong A_0*B_0,\qquad
 A_0=\frac{\pi_1(N_0)}{\langle\!\langle\im\pi_1(Y)\rangle\!\rangle},
 \quad
 B_0=\frac{\pi_1(C)}{\langle\!\langle\im\pi_1(Y)\rangle\!\rangle}.
\]
This use of van Kampen does not require either boundary homomorphism
to be injective: the boundary elements are trivial in the pushout, so
the pushout is the free product of the displayed quotients. These
quotients identify with the images of the two inclusion homomorphisms.
One-endedness forces $A_0=1$ or $B_0=1$.

Suppose $B_0=1$. Since $M$ is a $K(G,1)$, the inclusion $C\to M$ is
nullhomotopic. The homotopy extension property gives a map $f:M\to M$
homotopic to the identity and constant on $C$. Define
\[
 F:DN_0\longrightarrow M
\]
by $f|_{N_0}$ on the first copy and by the same constant on the second.
The definitions agree on $Y$. Moreover $\deg F=1$: collapsing $C$ in
$M$, or the second copy in $DN_0$, sends the fundamental class to the
same relative class $[N_0,Y]$ in $N_0/Y$. Since $f$ has degree one, so
does $F$.

This contradicts Chodosh--Li--Liokumovich
\cite[Corollary 3]{CLL}: no closed positive-scalar-curvature
four-manifold admits a nonzero-degree map to a closed aspherical
four-manifold. Thus $B_0\ne1$, so $A_0=1$. The inclusion of $N_0$
induces the trivial fundamental-group homomorphism and is consequently
nullhomotopic, again because the target is a $K(G,1)$.
\end{proof}

\begin{corollary}\label{cor:homology}
In Theorem~\ref{thm:domain}, if $Y$ is a spherical space form, then
\[
 H_2(N_0;\Z)=H_3(N_0;\Z)=0,\quad
 \chi(N_0)=1,\quad \tau(N_0)=0,
\]
and, for a finite abelian group $A$,
\[
 H_1(N_0;\Z)=A,\qquad H_1(Y;\Z)\cong A\oplus A.
\]
\end{corollary}

\begin{proof}
The inclusion lifts to an embedding in the contractible cover $U$.
For $C_U=\overline{U\setminus N_0}$, the integral Mayer--Vietoris
sequence and $H_2(Y;\Z)=0$ give
\[
 H_2(N_0;\Z)=0,\qquad
 H_1(Y;\Z)\cong H_1(N_0;\Z)\oplus H_1(C_U;\Z).
\]
Thus $A=H_1(N_0;\Z)$ is finite. Poincare--Lefschetz duality, together
with connected boundary and $H^1(N_0;\Z)=0$, gives $H_3(N_0;\Z)=0$.
The pair sequence, duality and the universal coefficient theorem give
\[
 0\longrightarrow\Ext(A,\Z)\longrightarrow H_1(Y;\Z)
   \longrightarrow A\longrightarrow0.
\]
The last map is projection to the first factor in the Mayer--Vietoris
splitting; its kernel is therefore isomorphic to $\Ext(A,\Z)\cong A$.
This proves the assertion. No compactification of $U$ is used.
\end{proof}

\section{Application to all scalar-nonnegative ALE cores}
Let $(N,g)$ be a smooth connected complete one-ended ALE four-manifold,
with end $S^3/\Gamma$. Assume differentiated positive-order ALE decay
sufficient for the usual Chern--Weil boundary limits. Write
\[
 W_\pm=\int_N|W^\pm|^2\,dV,
 \quad S=\int_Ns^2\,dV,
 \quad P=\int_N|\Ric^\circ|^2\,dV.
\]
The normalization is fixed by the characteristic identities below.

\begin{theorem}\label{thm:ale}
Suppose $s_g\ge0$ and a compact core of $N$ embeds smoothly into a
closed oriented aspherical four-manifold. If $\Gamma\ne1$, then
\[
 \boxed{P\ge4\pi^2+\frac{S}{12}.}
\]
There is no finite-fundamental-group hypothesis on $N$.
Every Ricci-flat ALE manifold satisfying the embedding hypothesis is flat.
\end{theorem}

\begin{proof}
Large ALE coordinate spheres have strictly positive outward mean
curvature. Apply Theorem~\ref{thm:domain} and
Corollary~\ref{cor:homology} to a sufficiently large truncation.
In particular $\chi(N)=1$, $\tau(N)=0$, and $\Gamma_{\rm ab}$ is a
direct double. The finite-group paragraph in the proof of
\cite[Theorem 5.2]{Hillman} implies
\[
 \Gamma=1,\quad I^*,\quad\text{or }Q(8k),\quad k\ge1.
\]
Only this group-theoretic consequence is used, not the stronger
classification of embeddings in a homology four-sphere.
The ALE Gauss--Bonnet and signature formulas are
\[
 8\pi^2(1-|\Gamma|^{-1})=W_++W_- -\tfrac12P+\tfrac1{24}S,
 \qquad 0=\frac{W_+-W_-}{12\pi^2}-\eta(S^3/\Gamma).
\]
Consequently
\begin{equation}\label{eq:eta}
 P\ge24\pi^2|\eta(S^3/\Gamma)|
       -16\pi^2(1-|\Gamma|^{-1})+\frac{S}{12}.
\end{equation}
For the quaternionic cases, Wright's spherical-space-form eta formula
\cite[Lemma 1]{Wright}, with $u=1$ and $v=2k$, gives
\[
 |\eta(S^3/Q(8k))|=\frac{2k}{3}+\frac1{12k}.
\]
The free orthogonal actions here are equivalent up to group automorphism
and orientation, which does not change the absolute eta invariant.
Thus
\[
 P\ge16\pi^2(k-1)+\frac{4\pi^2}{k}+\frac{S}{12}
   \ge4\pi^2+\frac{S}{12}.
\]
For $I^*$, $|\Gamma|=120$ and $|\eta|=361/180$, giving
$P\ge(484/15)\pi^2+S/12$. The eta value also follows from the
hyperkahler $E_8$ resolution, with $\chi=9$ and $\tau=-8$.
If $g$ is Ricci-flat, all nontrivial end groups are excluded. For the
trivial group the first characteristic identity then forces
$W_+=W_-=0$, hence flatness.
\end{proof}

The estimate is sharp. LeBrun's scalar-flat Kahler ALE metric on
$\mathcal O(-4)$ admits the free orientation-preserving antiholomorphic
isometric involution induced by
\[
 j(z_1,z_2)=(-\overline z_2,\overline z_1).
\]
Since $j^2=-\mathrm{Id}$ belongs to the end group $\Z_4$, this is an
involution on the ALE space. It is antipodal on the zero section and
free elsewhere. The quotient has end $S^3/Q_8$ and retracts onto
$\mathbb{RP}^2$. Its compact core is the normal disk bundle of a
standard embedded $\mathbb{RP}^2\subset S^4$, so embeds in $\R^4$ and
in every smooth four-manifold. On $\mathcal O(-m)$ the compactly
supported square is $c_1^2=-(m-2)^2/m$, and scalar-flat Kahler geometry
gives $P=-8\pi^2c_1^2$. Taking $m=4$ and the degree-two quotient gives
\[
 P\bigl(\mathcal O(-4)/\langle j\rangle\bigr)=4\pi^2.
\]
These are classical metric and embedding ingredients
\cite{LeBrun88,Hillman}; the quotient shows exactly what the theorem
does not exclude.

\begin{remark}
Scalar-flat Bach-flat ALE geometry alone does not imply finite
fundamental group. Blow up one point of the positive-scalar locally
conformally flat product $S^1\times S^3$ using the positive Green
function of its conformal Laplacian. The resulting complete metric is
scalar-flat, locally conformally flat, one-ended AE, and has fundamental
group $\Z$. Theorem~\ref{thm:domain} uses the ambient topology to
exclude an embedded core of this type, rather than assuming it away.
For $\Gamma=1$, the characteristic identities themselves only give
$P=2(W_++W_-)+S/12$; they provide no positive quantized bound for a
general nonflat core.
\end{remark}

\section{A global loss of compactness with no curvature bubbles}
One cannot infer attainment from the preceding ALE restrictions.
Here is an explicit comparison within the same unrestricted functional.

Take an orientable finite-volume real-hyperbolic four-manifold $V$ with
torus cusp sections. Its cusp metric has the form
\[
 dt^2+e^{-2t}h,
\]
where $h$ is flat on $T^3$. Double a compact truncation of $V$ along
all its cusp boundaries. The resulting closed manifold $M$ is
aspherical and has $\chi(M)=2\chi(V)>0$. It is a nontrivial pure
real-hyperbolic extended graph manifold, so it has no Einstein metric
by Di Cerbo \cite{DiCerbo}.

Choose a fixed smooth positive even function $a$ on $[-1,1]$, agreeing
with $e^{|x|}$ near the endpoints. At cusp depth $T$, join the two
hyperbolic sides with
\[
 g_T=dx^2+e^{-2T}a(x)^2h,\qquad -1\le x\le1.
\]
Outside this bridge retain the hyperbolic metric, truncating at the
matching depth. Thus $g_T$ is smooth on the same closed $M$ for every
large $T$.

Put $b=a'/a$. The radial and tangential Ricci eigenvalues on the
bridge are respectively
\[
 -3(b'+b^2),\qquad -b'-3b^2\quad\text{(multiplicity three)}.
\]
The trace-free eigenvalues are $-3b'/2,b'/2,b'/2,b'/2$, so
\[
 |\Ric^\circ|^2=3(b')^2,\qquad
 dV_{g_T}=e^{-3T}a^3\,dx\,dV_h.
\]
All exterior contributions to $P$ vanish, since the exterior is
hyperbolic. Consequently the complete, not merely local, energy is
\[
 \boxed{P(g_T)=3e^{-3T}
  \sum_{\mathrm{cusps}}\Vol_h(T^3)
  \int_{-1}^{1}a^3\bigl((a'/a)'\bigr)^2\,dx\longrightarrow0.}
\]
Hence $P_*(M)=0$, and the infimum is not attained: attainment would
give an Einstein metric. The sectional curvatures remain uniformly
bounded, because they are $-a''/a$ and $-(a'/a)^2$ on the bridge.
The missing compactness is the collapsing torus neck and unbounded
length, not curvature blow-up.

This is a direct consequence of classical cusp geometry and the cited
Einstein obstruction, not a claimed new example. It demonstrates a
precise limitation of a bubble-only existence argument, even when the
manifold is aspherical and has positive Euler characteristic.

\section{What this does and does not establish for the target}
The new connection is the use of scalar curvature on the double and
mapping degree to identify the topological side of an embedded ALE
core. It extends the previous energy theorem to every smooth
one-ended scalar-nonnegative ALE core, including cores whose
fundamental groups were not known to be finite.

This does not make Anderson's ideal attainment theorem
\cite[Theorem 4.10 and Remark 4.12]{Anderson} a smooth attainment
theorem. The quaternionic positive-energy core survives, the
trivial-end case has no bound of this form, intermediate ALE
orbifolds require their own embedding argument, and open or partially
collapsed limits remain possible. A topological embedding of the
quaternionic core inside a ball is not a controlled energy-decreasing
metric replacement: its ALE interface is $S^3/Q_8$, not the outer
$S^3$ of that ball.

There is also a useful numerical restriction on candidate selection.
Chern--Weil gives, on every closed oriented four-manifold,
\[
 P(g)\ge8\pi^2\bigl(3|\tau(M)|-2\chi(M)\bigr).
\]
If the strict Hitchin--Thorpe inequality is violated, the integer in
parentheses is at least one, so $P_*(M)\ge8\pi^2$. Such a candidate
cannot satisfy a proposed total-energy bound $P_*<4\pi^2$ for
excluding all nontrivial-end bubbles. In particular Capovilla's
aspherical candidates $\chi=\tau=n>0$ \cite{Capovilla} satisfy
$P_*\ge8\pi^2n$. Their positive lower bound distinguishes them from
the cusp doubles above, but does not itself prove attainment.

The unrestricted global existence assertion remains unproved.
The logical improvement is an enlarged global topological restriction
on all smooth ALE cores and a concrete account of why that restriction
alone cannot settle the minimization problem. No separate-local-minimum,
orbifold-minimum, or different-functional result is substituted for
the original target.

\begin{thebibliography}{10}
\bibitem{CLL} O. Chodosh, C. Li, Y. Liokumovich,
\emph{Classifying sufficiently connected PSC manifolds in 4 and 5 dimensions},
Geom. Topol. 27 (2023), 1635--1655.
\url{https://arxiv.org/abs/2105.07306}. Corollary 3.
\bibitem{Frenck} G. Frenck,
\emph{The doubling conjecture for positive scalar curvature} (2026).
\url{https://arxiv.org/abs/2604.12901}.
Introduction, recalling Gromov--Lawson's doubling implication.
\bibitem{Hillman} J. A. Hillman,
\emph{Locally flat embeddings of 3-manifolds in $S^4$}.
\url{https://arxiv.org/abs/2408.10535}.
Finite-group paragraph of the proof of Theorem 5.2.
\bibitem{Wright} E. P. Wright,
\emph{Quotients of gravitational instantons},
Ann. Global Anal. Geom. 41 (2012), 91--108.
\url{https://arxiv.org/abs/1102.2442}. Lemma 1.
\bibitem{LeBrun88} C. LeBrun,
\emph{Counter-examples to the generalized positive action conjecture},
Comm. Math. Phys. 118 (1988), 591--596.
\bibitem{LeBrun04} C. LeBrun,
\emph{Curvature functionals, optimal metrics, and the differential topology
of 4-manifolds}.
\url{https://arxiv.org/abs/math/0404251}.
\bibitem{DiCerbo} L. F. Di Cerbo,
\emph{Extended graph 4-manifolds, and Einstein metrics}.
\url{https://arxiv.org/abs/2106.13279}.
\bibitem{Anderson} M. T. Anderson,
\emph{Canonical metrics on 3-manifolds and 4-manifolds}.
\url{https://arxiv.org/abs/math/0511034}.
\bibitem{Capovilla} P. Capovilla,
\emph{Aspherical 4-manifolds with positive Euler characteristic and
their geography}.
\url{https://arxiv.org/abs/2511.15577}.
\end{thebibliography}
\end{document}
